% Figure: the two bending moments of an axisymmetric circular plate, the
% radial moment M_r and the circumferential moment M_theta, plotted along
% the radius for a clamped plate under uniform load. M_r is positive
% (sagging) at the centre and swings to a large negative (hogging) value at
% the clamped rim; M_theta stays positive across most of the radius. The
% rim is the governing section, which is why a brittle clamped plate cracks
% at the edge.
\begin{figure}[htbp]
\centering
\begin{tikzpicture}
\begin{axis}[
  width=11cm, height=6.5cm,
  xlabel={normalised radius $\ r/a$},
  ylabel={moment / $(q_0 a^2)$},
  xmin=0, xmax=1.0, ymin=-0.14, ymax=0.09,
  axis lines=left,
  grid=major, grid style={enggray!20},
  tick label style={font=\scriptsize},
  label style={font=\small},
  legend style={font=\scriptsize, at={(0.02,0.04)}, anchor=south west, draw=enggray!40},
]
% Clamped circular plate, uniform load, nu = 0.3.
% M_r/(q0 a^2)     = (1/16)[ (1+nu) - (3+nu) x^2 ]
% M_theta/(q0 a^2) = (1/16)[ (1+nu) - (1+3nu) x^2 ]
% radial moment
\addplot[engblue, very thick, domain=0:1, samples=80]
  {(1/16)*((1+0.3) - (3+0.3)*x*x)};
\addlegendentry{radial $M_r$}
% circumferential moment
\addplot[engred, very thick, domain=0:1, samples=80]
  {(1/16)*((1+0.3) - (1+3*0.3)*x*x)};
\addlegendentry{hoop $M_\theta$}
% zero line
\addplot[enggray, dashed, domain=0:1, samples=2] {0};
% mark the clamped-rim radial moment -q0 a^2 / 8 = -0.125
\addplot[engamber, only marks, mark=*, mark size=1.6pt] coordinates {(1.0,-0.125)};
\node[engamber, font=\scriptsize, anchor=south east] at (0.98,-0.118)
  {rim moment $-\tfrac{1}{8}q_0 a^2$};
% mark the centre value
\node[engblue, font=\scriptsize, anchor=south west] at (0.03,0.082)
  {centre $\tfrac{1+\nu}{16}q_0 a^2$};
\end{axis}
\end{tikzpicture}
\caption[Moments in an axisymmetric clamped circular plate]{The two bending
moments of a clamped circular plate under uniform load $q_0$, normalised by
$q_0 a^2$ and plotted along the radius ($\nu=0.3$). The radial moment $M_r$
(blue) is positive and sagging at the centre and swings to its largest
magnitude, a hogging $-q_0 a^2/8$, at the clamped rim (amber marker). The
circumferential moment $M_\theta$ (red) stays positive across the plate and
is smaller in magnitude. The governing section is the rim, not the centre,
which is why a brittle clamped plate cracks first at its edge, while a
simply supported plate, whose rim moment is zero, cracks at its centre. The
two curves are read straight from the axisymmetric plate equation, the
single-variable reduction of the biharmonic.}
\label{fig:vol03ch09:axisymmetric-plate}
\end{figure}
